\section{Restriction from a torus}\label{sec:torus}

The subgroup restriction problem of \cite[Section 3.4]{GCT6} allows $H$ to be a
torus and $\rho:H\to G$ any explicit polynomial homomorphism
\cite[Section 3.4.1]{GCT6}; the rank of the input is the sum of the ranks of
$H$ and $G$ and the heights of the specified weights, and it does not see the
exponents of $\rho$. For $H=\Gm$ the multiplicity of the character $t\mapsto
t^{k}$ in the restriction of $\Sym^n\C^N$ is the number of monomials of degree
$n$ whose weight is $k$.

\begin{proposition}\label{prop:torusmodular}
For $N\ge2$ let $\rho_N:\Gm\to GL_2$, $\rho_N(t)=\diag(t^{N+1},t)$, let $\pi=(1)$
(the representation $\C^2$, stretched to $\Sym^n\C^2$) and $\lambda$ the
character $t\mapsto t^2$ (stretched to $t^{2n}$). Then
\[
 f_N(n)=\dim\Hom_{\Gm}\bigl(t^{2n},\Sym^n(\C^2)|_{\rho_N}\bigr)=\mathbf1_{N\mid n},
 \qquad \sum_{n\ge0}f_N(n)T^n=\frac1{1-T^N},
\]
and every expression of this series as $P(T)/\prod_i(1-T^{a_i})^{b_i}$ has some
$a_i$ divisible by $N$, while $\rank(x)=1+2+1+1=5$. Hypothesis 3.3.4 of
\cite{GCT6} fails for subgroup restriction as stated.
\end{proposition}

\begin{proof}
The monomial $y_1^ay_2^{n-a}$ has weight $(N+1)a+(n-a)=n+Na$, which equals
$2n$ exactly when $Na=n$; there is one such $a\in[0,n]$ if $N\mid n$ and none
otherwise. The series has a pole at every primitive $N$-th root of unity, and a
denominator must vanish there.
\end{proof}

Here $p(f_N)=s(f_N)=0$, so this example does not touch PH2 or SH.

\begin{proposition}\label{prop:torusindices}
For odd $N\ge3$ let $\sigma_N:\Gm\to GL_3$, $\sigma_N(t)=\diag(t,t^3,t^{N+1})$,
with $\pi=(1)$ and $\lambda=t^2$ as before. Then
\[
 g_N(n)=\#\{(b,c)\in\Z_{\ge0}^2:2b+Nc=n\},\qquad
 \sum_{n\ge0}g_N(n)T^n=\frac1{(1-T^2)(1-T^N)},
\]
and $p(g_N)=s(g_N)=N-2$, while $\rank(x)=1+3+1+1=6$. Hypotheses 3.3.2(a) and
3.3.3 of \cite{GCT6} fail for subgroup restriction as stated.
\end{proposition}

\begin{proof}
A monomial $y_1^ay_2^by_3^c$ of degree $n$ has weight $a+3b+(N+1)c=n+2b+Nc$,
which equals $2n$ exactly when $2b+Nc=n$, and then $a=n-b-c\ge0$ automatically.
This gives $g_N$ and its series. As $N$ is odd, $n-Nc$ is even exactly when
$c\equiv n\pmod2$, so $g_N(n)$ is the number of $c\equiv n\pmod2$ with
$0\le c\le n/N$. The period is $2N$. For $n=2Nq+r$, $0\le r<2N$, one finds
$g_N(n)=q+1$ if $r$ is even, and $g_N(n)=q+\mathbf1_{r\ge N}$ if $r$ is odd. So
the constituents are $g_r(X)=(X-r)/(2N)+\varepsilon_r$, with $\varepsilon_r=1$
for even $r$ and $\varepsilon_r=\mathbf1_{r\ge N}$ for odd $r$.

The constituents of $n\mapsto g_N(n+p)$ are $X\mapsto g_r(X+p)$, $0\le r<2N$,
with constant terms $(p-r)/(2N)+\varepsilon_r$ and slope $1/(2N)$. For odd
$r<N$ the constant term is $(p-r)/(2N)$, nonnegative only if $p\ge r$; the
largest such $r$ is $N-2$. For $p=N-2$ the other constant terms are positive:
$(N-2-r)/(2N)+1\ge1/2$ for even $r\le2N-2$, and $\ge(N-1)/(2N)$ for odd
$r\ge N$. So $p(g_N)=N-2$. A constituent $X\mapsto g_r(X+s)$ is not
identically zero and is increasing, so it is positive at every integer $X\ge1$
exactly when $(1+s-r)/(2N)+\varepsilon_r>0$; for odd $r<N$ this means $s\ge r$,
and for $s=N-3$ it fails at $r=N-2$; for $s=N-2$ all hold. So $s(g_N)=N-2$.
\end{proof}

Neither example touches the density clause 3.3.2(b). They show that, for the
general subgroup restriction problem, the polynomial bounds must also depend on
the bit length of the embedding; with that dependence these two examples are
harmless. The Kronecker and plethysm families of
Sections~\ref{sec:kronecker}--\ref{sec:indices} have no such parameter: their
embeddings are fixed by $n$.

\section{Earlier work and the scope of the results}\label{sec:literature}

\paragraph{What we searched.}
On 25 November 2026 we searched the web and arXiv for counterexamples to
Mulmuley's positivity and saturation hypotheses (queries on PH2, PH3, SH,
modular index, stretching quasi-polynomials, Kronecker and plethysm
coefficients), and the Hexagon repository (\url{https://hexagonmath.org}) for
``Mulmuley'', ``plethysm'', ``Kronecker'', ``geometric complexity'',
``positivity hypothesis'', ``saturation hypothesis'' and ``modular index''. The
Hexagon searches returned no work on these hypotheses (the matches for
``plethysm'' and ``Kronecker'' concern other questions). This is a bounded
search, and we make no claim of priority beyond it.

\paragraph{Briand, Orellana and Rosas.}
In \cite{BOR}, Briand, Orellana and Rosas showed that the strict saturation
hypothesis of the first versions of GCT~VI (``strong SH'') fails for Kronecker
coefficients: for $\lambda=(k,k)$, $\mu=(k+1,k-1)$, $\nu=(2k-2i-2j,2i,2j)$
with $i>j>0$ and $k>2i+j$, the stretched coefficient is $N/2+1$ for even $N$
and $-1/2$ for odd $N$ \cite[Theorem 2.4, p.~7]{BOR}; they also gave
counterexamples to strong PH2 \cite[Section 2.4, p.~8]{BOR}. In their examples
the positivity index is at most one \cite[Section 2.5, p.~8]{BOR}. Mulmuley's
appendix \cite[pp.~18--20]{BOR} then relaxed SH, PH2 and PH3 to polynomial
bounds in the largest height. Our Theorem~\ref{thm:kronecker} refutes the
relaxed PH3 of that appendix; the relaxed SH(a) and PH2 for Kronecker
coefficients remain open as far as this paper is concerned, while their
analogues for plethysm fail by Theorem~\ref{thm:indices}. GCT~VI v4 (22 January
2009), the newest version on arXiv, states the relaxed hypotheses in the form
quoted in Section~\ref{sec:defs}.

\paragraph{Kahle and Micha\l{}ek.}
Kahle and Micha\l{}ek~\cite{KM} showed, on the examples $S^3(S^k)$ and
$S^k(S^3)$, that the quasi-polynomials of plethysm need not count lattice points
of inhomogeneous polytopes whose dimension equals their degree. This concerns
the polytope hypothesis PH1 of \cite{GCT6}, not the index hypotheses treated
here.

\paragraph{What remains open.}
(1) The exponent $r$ in Theorem~\ref{thm:kronecker}: whether $K_n$ always has
its pole at the primitive $2^{n-1}$-th roots themselves. (2) The positivity and
saturation indices of the Kronecker family $k_n$. (3) The density clauses
1.6.5(b) and 3.3.2(b). (4) Whether a weakened PH3, with the modular index
replaced by a quantity insensitive to a single pole of exponential order, could
serve the purposes of \cite{GCT6}; the poles found here are simple, isolated in
the two-adic direction, and come from one fixed curve.

\section{Who proved and checked what}\label{sec:provenance}

The Kronecker results (Theorems~\ref{thm:K023}, \ref{thm:K024},
\ref{thm:transfer}, \ref{thm:kronecker}, Lemma~\ref{lem:liezero}) were found
and proved by a run of GPT-6 Astra (OpenAI, ChatGPT Pro mode) on 21--22
November 2026; the support theorem (Section~\ref{sec:support}),
Proposition~\ref{prop:fixedcomponentpole} and the plethysm results
(Sections~\ref{sec:plethysm}--\ref{sec:indices}) and the torus examples by an
earlier run of the same model on 20--21 November. Claude Opus~5.5 at effort max
read every proof of the Kronecker chain, of
Proposition~\ref{prop:fixedcomponentpole} and of Section~\ref{sec:support} and
reran the computations with its own code (22 November); Claude Opus~5.5 at
effort high read the proofs of Theorem~\ref{thm:plethysmPH3}, of the
canonical-algebra identification and of Section~\ref{sec:indices} and recomputed
their numerical claims in small dimensions (21 November). In writing this text
the latter added: the Kempf--Ness proof of closedness for the chain cubic
(replacing a citation), the grading stabilizer of $p_n$ for all $n$ and
Lemma~\ref{lem:orbitclosure} (which give $e_Y=1$ also when $3\mid n$), the
direct count of fixed weights, the route to \eqref{eq:amplitude} through
Proposition~\ref{prop:product}, the Fourier normalization, and the proofs of
Propositions~\ref{prop:torusmodular} and~\ref{prop:torusindices}; it checked
these, and the Hessian locus of Proposition~\ref{prop:chaincubic}(2), by
computation. No proof in this paper has been checked by a human.

\begin{thebibliography}{99}

\bibitem{BOR}
E.~Briand, R.~Orellana, M.~Rosas, \emph{Reduced Kronecker coefficients and
counter-examples to Mulmuley's strong saturation conjecture SH}, with an
appendix by K.~Mulmuley, Comput. Complexity \textbf{18} (2009);
arXiv:0810.3163v3 (14 May 2009). Page numbers refer to the arXiv version.

\bibitem{KM}
T.~Kahle, M.~Micha\l{}ek, \emph{Obstructions to combinatorial formulas for
plethysm}, Electron. J. Combin. \textbf{25}(1) (2018), \#P1.41;
arXiv:1507.07131v3.

\bibitem{Losev}
I.~Losev, \emph{The Kempf--Ness theorem and invariant theory},
arXiv:math/0605756.

\bibitem{Luna}
D.~Luna, \emph{Slices \'etales}, Bull. Soc. Math. France, M\'emoire \textbf{33}
(1973), 81--105.

\bibitem{Macdonald}
I.~G.~Macdonald, \emph{Symmetric functions and Hall polynomials}, 2nd ed.,
Oxford University Press, 1995.

\bibitem{GCT6}
K.~D.~Mulmuley, \emph{Geometric complexity theory VI: the flip via saturated
and positive integer programming in representation theory and algebraic
geometry}, arXiv:0704.0229v4 (22 January 2009). Page numbers refer to this
version.

\bibitem{Stacks}
The Stacks Project Authors, \emph{The Stacks Project},
\url{https://stacks.math.columbia.edu}, Section 15.29 (Tag 0621) and
Lemma 15.31.2 (Tag 062F).

\end{thebibliography}

\end{document}
